% ============================================================
% 02_related.tex
% ============================================================

\section{Related Work}
\label{sec:related}

\paragraph{Mode effects in survey measurement.}
Survey methodology settled long ago that the instrument is part of the measurement. Answers to the same question differ by mode of administration because each mode changes the cognitive and social conditions of answering \citep{tourangeau2000}. \citet{klausch2013} showed that telephone, face-to-face, and web modes yield non-equivalent measurement of identical attitude scales, and mixed-mode designs now model mode explicitly instead of ignoring it as a nuisance \citep{hox2015}. The Total Survey Error framework gives these effects a named place in the error budget \citep{groves2010}. \citet{sen2021} transported this thinking to digital traces and explicitly identified platform-affordance error as the trace analogue of mode error, but the framework names the error without supplying a test for it. Our premise is simply that platform should receive the same formal treatment for machine-labelled trace data that mode receives for surveys.

\paragraph{Measurement for machine-labelled text.}
A second literature treats classifiers as measurement instruments. \citet{jacobs2021} argue that computational pipelines operationalise unobservable constructs through measurement models whose validity must be tested, and \citet{halterman2025} extend this framing to large language models. Empirical work in this vein audits the instrument itself. \citet{pellert2022} validate sentiment macroscopes against self-report. \citet{atreja2025} show that model-generated labels shift with prompt design in unpredictable ways, and \citet{baumann2025} show that defensible configuration choices alone can flip downstream conclusions. \citet{egami2023} correct downstream estimates for classifier error with design-based supervised learning, but the correction targets label error against a gold standard and leaves platform-specific distortion untouched. In all of this work the classifier is scrutinised; the platform it reads from is not.

\paragraph{Cross-platform studies and assumed equivalence.}
Multi-platform designs have become standard, and the post-API turn has made heterogeneous data regimes the norm \citep{davidson2023, goanta2026}. \citet{overgaard2024} apply a single BERT classifier across Facebook and Twitter and compare the resulting scores directly; \citet{wei2026} compare Douyin and TikTok in the same way. There is growing evidence that this is risky. \citet{greene2025} find that the same actors are read differently across platforms, and \citet{nogara2025} show that Perspective API scores German-language text as more toxic, a measurement artefact that masquerades as a substantive difference. The field's classic cautions concern populations and samples: who is on the platform \citep{tufekci2014}, how the stream is sampled \citep{morstatter2013}, whether models generalise beyond their training population \citep{cohen2013}, and whether text sentiment tracks opinion at all \citep{oconnor2010}. The complementary question, whether the same score means the same thing on Reddit as on YouTube, is routinely assumed and almost never tested.

\paragraph{The psychometric toolkit and nearest misses.}
The tools for this question already exist in psychometrics. Generalizability theory partitions score variance among measurement facets \citep{cronbach1972, brennan2001}; measurement invariance testing asks whether indicators relate to a construct identically across groups \citep{meredith1993, putnick2016}; linking places scores from different instruments on a common scale \citep{kolen2014}; and invariance results are read as diagnostic evidence rather than permission to compare \citep{robitzsch2023}. Several recent studies come close to combining this toolkit with platform data without quite doing so. \citet{boyd2024} do test cross-platform invariance, though of survey self-reports rather than machine labels. Faceted Rasch measurement has reached annotation through \citet{kennedy2020} and \citet{sachdeva2022}, with annotators as the modelled facet and platform left out of the model. In \citet{pellert2022} validation is by correlation alone, with no invariance structure. The multitrait-multimethod models of \citet{cernat2025} compare survey with digital-trace measures rather than platform with platform. \citet{han2025} document platform-dependent selection and measurement bias, but in online review ratings rather than machine-labelled text. \citet{pokropek2026} fit factor models to text indicators while grouping by time within a single platform. And generalizability designs have reached LLM annotation, with prompt and model choice as variance facets, without yet reaching platform \citep{camuffo2026}.

To our knowledge, no study has applied the psychometric equivalence toolkit of generalizability, invariance, and linking to machine-classified policy sentiment with platform as the grouping facet, and this paper specifies that test.
